\emph{Notation.} The frame of \claimref{synthesis.positive-conformal-energy-gives-positive-null-translations}: $e_{0}^{2} = \nu^{2} = 1$, $e_{1}^{2} = e_{2}^{2} = e_{3}^{2} = e_{-}^{2} = -1$, $I_{4} = e_{0}e_{1}e_{2}e_{3}$, $E = \nu e_{-}$, $I_{6} = I_{4}E$ (the conventions' $I_{6}$), $Z = e_{0}\nu$. $L_{a}\psi = a\psi$, $R_{b}\psi = \psi b$ on $\Cl(2,4)$; the conformal spinors' unit is $R_{I_{6}}$, which commutes with every $L_{a}$ and has $R_{I_{6}}^{2} = R_{I_{6}^{2}} = -1$. The planes are $(x_{1}, y_{1}) = (e_{0}, \nu)$, $(x_{2}, y_{2}) = (e_{1}, e_{2})$, $(x_{3}, y_{3}) = (e_{3}, e_{-})$, with $x_{a}^{2} = y_{a}^{2}$; write $(s_{1}, s_{2}, s_{3}) = (s_{0}, s_{1}, s_{2})$ for the three signs, $s_{a}$ attached to the $a$-th plane.

\emph{Step 1 (the flag on the module).} An orthogonal map $J$ of a plane $(x, y)$ with $x^{2} = y^{2}$ and $J^{2} = -1$ is a quarter turn, $Jx = \pm y$; so the orthogonal complex structures of $\R^{2,4}$ preserving the three planes are exactly the eight $J_{s}$, $J_{s}x_{a} = s_{a}y_{a}$, $J_{s}y_{a} = -s_{a}x_{a}$. With the unit $R_{I_{6}}$ in place of the scalar unit, the isotropic vector $x_{a} - s_{a}y_{a}$ of the flag $W_{s}$ (the eigenspace of $J_{s}$ for the unit) acts on the module as $w_{a} = L_{x_{a}} - s_{a}L_{y_{a}}R_{I_{6}}$. Since $R_{I_{6}}$ commutes with the left action, $w_{a}^{2} = L_{x_{a}^{2}} + L_{y_{a}^{2}}R_{I_{6}}^{2} - s_{a}L_{x_{a}y_{a} + y_{a}x_{a}}R_{I_{6}} = x_{a}^{2} - y_{a}^{2} = 0$, and for $a \ne b$ every term of $w_{a}w_{b} + w_{b}w_{a}$ contains $L_{u v + v u}$ with $u \perp v$, so it vanishes: $W_{s}$ is isotropic, the module form of $W_{s}\psi = 0$. Moreover $w_{a}\psi = 0$ says $x_{a}\psi = (J_{s}x_{a})\psi I_{6}$, and right multiplication by $-I_{6}$ turns it into $y_{a}\psi = -s_{a}x_{a}\psi I_{6} = (J_{s}y_{a})\psi I_{6}$; by linearity the pure spinors are the $\psi$ with $v\psi = (J_{s}v)\psi I_{6}$ for every vector $v$. Their space $K_{s}$ is stable under $R_{c}$ for even $c$ ($c$ commutes with $I_{6}$); the exact computation gives $\dim K_{s} = 8$ for each of the eight flags, generated by any one nonzero element under this right action, and in $\CS = \Cl^{+}(2,4)P$ (respectively its odd counterpart), a module over $\mathrm{span}\{1, I_{6}\}$ through $R_{I_{6}}$, the pure spinor is unique up to $\mathrm{span}\{1, I_{6}\}$, as Proposition \claimref{dark-matter.two-families-are-two-chiralities}(i) states for the complexified module.

\emph{Step 2 (the conformal energy).} From $e_{0}\psi = s_{0}\nu\psi I_{6}$, right multiplication by $-I_{6}$ gives $\nu\psi = -s_{0}e_{0}\psi I_{6}$, so $Z\psi = e_{0}\nu\psi = -s_{0}e_{0}^{2}\psi I_{6} = -s_{0}\psi I_{6}$. Hence $(Z\psi)I_{6} = s_{0}\psi$, and for the positive form $\ip{x}{y} = \br{x^{\dagger}y}_{0}$ ($\dagger$ the reversion with a sign for each generator of square $-1$, so that frame vectors of square $+1$ are self-adjoint and those of square $-1$ anti-self-adjoint, as in \claimref{synthesis.attachment-conformal-energy-is-generated-by-z}), $\ip{\psi}{(Z\psi)I_{6}} = s_{0}\ip{\psi}{\psi}$: the real part of the twistor form $h(\psi, \psi) = \ip{\psi}{\mathrm{unit}\cdot Z\psi}$ has the sign $s_{0}$. In the normalisation of that result's clause (i), $Z$ acts as minus the unit on the positive plane $T_{+}$. That unit is a scalar of a complex module there; identifying it with $+R_{I_{6}}$ is a convention, since the commutant fixes the unit only up to sign (\claimref{field-theory.orientation-relocated}(3)). With it, $Z = -s_{0}R_{I_{6}}$ on the vacuum, which lies in $T_{+}$ exactly when $s_{0} = +1$; with $-R_{I_{6}}$, $T_{+}$ and $T_{-}$ exchange. Here $Z = P_{e_{0}} + e_{0}\nO$ with $P_{a} = \tfrac12 a\nInf$, the conformal energy of \claimref{synthesis.positive-conformal-energy-gives-positive-null-translations}(i), whose clause (ii) ties the sign of $Z$ in a unitary representation to the sign of every future null translation.

\emph{Step 3 (the family).} As in Step 2, $e_{1}e_{2}\psi = -s_{1}e_{1}^{2}\psi I_{6} = s_{1}\psi I_{6}$ and $e_{3}e_{-}\psi = s_{2}\psi I_{6}$. Moving $\nu$ to the left past $e_{3}$, $e_{2}$, $e_{1}$, $I_{6} = e_{0}e_{1}e_{2}e_{3}\nu e_{-} = -(e_{0}\nu)(e_{1}e_{2})(e_{3}e_{-})$; the three bivectors commute with each other and with $R_{I_{6}}$, so $I_{6}\psi = -(-s_{0})(s_{1})(s_{2})\psi I_{6}^{3} = -s_{0}s_{1}s_{2}\,\psi I_{6}$. Moving $\nu$ the same way, $e_{0}\wedge J_{s}e_{0}\wedge e_{1}\wedge J_{s}e_{1}\wedge e_{3}\wedge J_{s}e_{3} = s_{0}s_{1}s_{2}\,e_{0}\nu e_{1}e_{2}e_{3}e_{-} = -s_{0}s_{1}s_{2}I_{6}$, so $I_{6}\psi = \omega(J_{s})\psi I_{6}$ with $\omega(J_{s})$ the orientation of $J_{s}$ relative to $I_{6}$ (the standard statement that the chirality of a pure spinor is the orientation of its complex structure; Cartan, Chevalley, cited from memory). The operators $w_{a}$ are odd ($L_{v}$ odd, $R_{I_{6}}$ even), so $K_{s}$ is the sum of its even and odd parts; $I_{6}$ commutes with even and anticommutes with odd elements, and $R_{I_{6}}$ is invertible, so $K_{s}$ is even when $\omega = 1$ and odd when $\omega = -1$. The even part is the half $\CS = \Cl^{+}(2,4)P$ and the odd part $\CSm$ (conventions), and the chirality of the pure spinor is the family of Proposition \claimref{dark-matter.two-families-are-two-chiralities}(ii) (Definition \claimref{obj.dark-matter.null-blades-fock-ideal-chirality}, whose overall sign is a convention, only ratios being used). Reversing all three signs reverses $\omega$: the conjugate flag lies in the other family, as Definition \claimref{obj.dark-matter.witt-flag}, the same Proposition (iii) ($m = 3$ odd) and \claimref{synthesis.cpt-conjugation-pseudoscalar-rotor-times-complex-conjugation}(2)-(3) state. With the opposite orientation of $I_{6}$ and the same labels $s$, the family reads $+s_{0}s_{1}s_{2}$ and Step 2 is unchanged. Finally $e_{0}\wedge J_{s}e_{0} = s_{0}e_{0}\nu = s_{0}Z$ and $e_{1}\wedge J_{s}e_{1}\wedge e_{3}\wedge J_{s}e_{3} = s_{1}s_{2}\,e_{1}e_{2}e_{3}e_{-}$: the orientation of $J_{s}$ on the plane of $Z$ is $s_{0}$ and on the four-space orthogonal to $Z$ it is $s_{1}s_{2}$.

\emph{Step 4 (conjugate pairs).} The energy sign is $s_{0}$ and the family is $\omega = -s_{0}\cdot s_{1}s_{2}$. As $s$ runs over the eight flags, $(s_{0}, \omega)$ takes each of its four values exactly twice, so as functions on the set of eight flags neither determines the other. The orientation bit of Premise \claimref{dark-matter.premise-flag-as-datum} is the sign of $J$ against $-J$: its two values are $W_{s}$ and $W_{-s}$, $J_{-s} = -J_{s}$. Under $s \mapsto -s$, $s_{0}$ and $\omega$ both change sign and $s_{1}s_{2}$ does not, so $s_{0}\,\omega = -s_{1}s_{2}$ is the same on both members of every conjugate pair: relative to the stored bit, the energy sign is the family times the fixed sign $-s_{1}s_{2}$.

\emph{Step 5 (reading; argued).} Definition \claimref{obj.dark-matter.witt-flag} makes the flag the annihilator of a Fock vacuum, and Premise \claimref{dark-matter.premise-flag-as-datum} places the sign of $J$ relative to the Fock unit in the configuration. If the orientation datum of Premise \claimref{synthesis.n-plus-null-ground-state} is read through the conformal energy of the flag's vacuum (positive conformal energy giving positive null translations, \claimref{synthesis.positive-conformal-energy-gives-positive-null-translations}(ii); the attachment's unit taken as $+R_{I_{6}}$), then by Step 4 it is the stored flag bit up to the sign $s_{1}s_{2}$ of the descent's complex structure on the four-space orthogonal to $Z$, which is kinematic once $J$ is fixed up to $J \mapsto -J$. This agrees with \claimref{one-relation.arrow-frequency-are-one-bit-level} and Fact \claimref{one-relation.imported-flag-arrow-superselection}, by which the arrow and the frequency sign are one bit. Whether the descent's flag is one of these eight $Z$-adapted flags, and what fixes $s_{1}s_{2}$, is open. Steps 1 to 4 are exact and complete for the eight flags; the checks verify them exactly, the conjugate-pair relation included.
